% Unified manuscript insert: two recursive retractions and structural rescue

\subsection{Constraint-aware recursive retractions}
\label{sec:constraint-retractions}

Let $Y=(Y_1,\ldots,Y_B)$ denote a sibling Monte Carlo batch of recursive
MLP child-state estimates, and let $\mathcal C(t,x)$ be a PDE-certified
feasible set for the exact state. We apply a retraction
\[
  \mathcal R_B:Y\longmapsto \mathcal R_B(Y)\in\mathcal C(t,x)^B
\]
\emph{before} the corrected child states are reused by the nonlinear
generator $F$. We require $\mathcal R_B(Y)=Y$ for an already feasible
batch. This formulation separates the structural requirement of returning to
the certified set from the statistical geometry used to perform the
correction.

\paragraph{Samplewise IR-MLP.}
For a closed convex set,
\[
 [\mathcal R_{\rm sep}(Y)]_i=\Pi_{\mathcal C}(Y_i).
\]
Equivalently,
\[
 \mathcal R_{\rm sep}(Y)
 \in\arg\min_{Z\in\mathcal C^B}\sum_{i=1}^B\|Z_i-Y_i\|_2^2.
\]
This retraction is local and nonexpansive.

\paragraph{Batchwise Gauge IR.}
For a star-shaped or gauge feasible set
\[
 \mathcal C=\{y:\gamma_{\mathcal C}(y)\le1\},
\]
define
\[
 \Gamma(Y)=\max_{1\le i\le B}\gamma_{\mathcal C}(Y_i),
 \qquad
 \alpha_B(Y)=\min\{1,\Gamma(Y)^{-1}\},
\]
and use
\[
 [\mathcal R_{\rm batch}(Y)]_i=\alpha_B(Y)Y_i.
\]
Thus all sibling samples use the largest common feasible scale. The method
preserves their relative directions and ratios but couples the batch and
generally changes its empirical mean. Samplewise and batchwise IR therefore
share the same certified product set, insertion point, and fixed-point target;
they differ in the geometry of the retraction.

\subsection{A retraction-invariant counterexample rescue}
\label{sec:retraction-rescue}

For the counterexample let
\[
 S_d=\operatorname{span}(e_1),
 \qquad
 f_d(z)=\|z_{2:d}\|_2.
\]
The control certificate gives $\nabla u_d(t,x)\in S_d$, and
$f_d(z)=0$ for every $z\in S_d$.

\begin{proposition}[Structural recursive rescue]
\label{prop:structural-retraction-rescue}
Consider any recursive correction, possibly coupled across a sibling Monte
Carlo batch, such that every corrected gradient returned to a parent generator
evaluation belongs to $S_d$. Then every nonlinear full-history Picard
correction vanishes pathwise. Consequently the high-dimensional nonlinear
recursion reduces to terminal Monte Carlo on the active first coordinate.
\end{proposition}

The proposition is deliberately independent of Euclidean projection.
Samplewise IR uses $z_i\mapsto Q_dz_i$, where $Q_d$ is the orthogonal
projection onto $S_d$. A structurally compatible batchwise realization is
\[
 z_i\mapsto \alpha_BQ_dz_i,
 \qquad 0\le\alpha_B\le1,
\]
which has the same pathwise nonlinear collapse. In contrast, pure common
scaling $z_i\mapsto\alpha_Bz_i$ does not exact-rescue this example because
it need not remove the inactive coordinates.

For the samplewise choice at $(t,x)=(0,0)$, with $N=m^n$ terminal samples,
\[
 \mathbb E\|\widehat Y^{\rm IR}-Y^*\|_2^2
 =\frac{4-2/\pi}{N},
\]
independent of the ambient dimension. The structural theorem therefore
identifies the essential property behind the rescue: recursive return to a
generator-compatible certified set, rather than a particular choice of
nearest-point projector.

\paragraph{Finite-sample tradeoff.}
The two retractions need not have the same finite-sample error away from the
exact-rescue setting. In paired low-budget HJB experiments, Batch-IR reduces
relative $L_2$ error from $0.2412$ to $0.0966$ at $d=20$ and from $0.1993$
to $0.0833$ at $d=40$ compared with Samplewise IR. In the 100-dimensional
nonlinear-funding benchmark, the corresponding MAEs are $0.5102$ versus
$0.4807$ at $M=10$, $0.2965$ versus $0.2539$ at $M=20$, and $0.1993$
versus $0.1486$ at $M=40$. We therefore treat the two algorithms as
complementary realizations of the same constraint-aware principle rather than
claiming that Euclidean projection is universally optimal.
